\documentclass{ximera}

\begin{document}

    \title{Lecture 13}
    
    \begin{abstract}  
    Outline: \\
    - Closed sets \\
    - Complements
\end{abstract}

\maketitle

Recall:

\begin{definition}
    A set $F\subseteq\mathbb{R}$ is closed if it contains all its limit points.
\end{definition}



\begin{definition}
    Let $A\subseteq \mathbb{R}$, and let $L$ be the set of limit points of $A$.  The closure of $A$ is the set $\overline{A}:= A\cup L$.
\end{definition}

\begin{theorem}
    For any $A\subseteq \mathbb{R}$, the closure of $A$ is a closed set and the smallest closed set containing $A$.
\end{theorem}
\begin{proof}
    Let $L$ be the set of limit points of $A$.  By definition, $L \subseteq \overline{A}$, but we need to show that $\overline{A}$ has not generated any new limit points.  Let $x$ be a limit point of $\overline{A}$.  Then, for all $\epsilon > 0$,  $V_\epsilon(x) \cap \overline{A} \neq \emptyset$.  Note that, using DeMorgan's laws, $V_\epsilon(x) \cap (A\cup L) = (V_\epsilon(x) \cap A) \cup (V_\epsilon(x) \cap L)$ .  If $x \in L$, then we are done.  If $x \notin L$, i.e. if it is an isolated point, then there exists $\epsilon_0$ such that $V_{\epsilon_0}(x) \cap A = \emptyset$; but then this would also imply that $V_{\epsilon_0}(x) \cap \overline{A} = \emptyset$, a contradiction. 

    Now, we need to show that $\overline{A}$ is the smallest closed set containing $A$.  This is trivial, as any closed set containing $A$ must also contain $L$, i.e. $A\cup L \subset F$ for any closed set $F$ containing $A$.  But $A \cup L = \overline{A}$ by definition, and the theorem is proved.
\end{proof}

\begin{exercise}
    Come up with an example of a set that is 
    \begin{enumerate}
        \item open
        \item closed
        \item contains only isolated points
        \item contains only limit points
        \item is neither closed nor open, but does not only consist of isolated points.
    \end{enumerate}
\end{exercise}

    \begin{exercise}
        Is it possible for a set in $\mathbb{R}$ to be closed and unbounded?
    \end{exercise}

    Recall the definition of complement of a set:
    \begin{definition}
        Let $A \subseteq \mathbb{R}$.  The complement of $A$ is defined to be 
        \[A^C = \{x \in \mathbb{R}: x \notin A\}.\]
    \end{definition}

    \begin{exercise}
        Find the complement of the sets you came up with in Exercise 3.  Determine whether or not they are open, closed, or neither.   What do you conjecture about the nature of the complement of open sets? closed sets?
    \end{exercise}
    
\end{document}